\documentclass[11pt]{article}
\usepackage{preprint}

\title{The Construction of Ethnicity, Race, and Nation}
\author{Jiadi He\email{jh2288@georgetown.edu}}
\shorttitle{The Construction of Ethnicity, Race, and Nation}
\keywords{identity; political action; social construction}

\begin{document}
\maketitle
\keywordsline

\section{Introduction}

It is an undeniable fact that the human population is diverse. In different places across the world, humans are generally the same, but they also exhibit different physical attributes (like skin color, and hair color and texture), and speak different languages. These elements have been used as a basis for the concepts of ethnicity, race, and nation. According to a significant body of literature, these concepts (of ethnicity, race, and nation) are not intrinsically true, rather they are social/political constructions that have happened over time, aided by widely shared narrative.\footnote{Brubaker, Rogers. Ethnicity without Groups (Cambridge: Harvard University Press, 2004) 9} Such notions imply a time when these distinctions of human populations were not in existence. One concept that has drawn much attention is race, as examined by Yudell and other scholars.\footnote{Yudell, Michael, Dorothy Roberts, Rob DeSalle, \& Sarah Tishkoff. “Taking Race Out of Human Genetics.” \textit{Science and Society vol.351}, No.6273, 2016, pp.351-353} In a nutshell, these concepts are no more than make-believe. In that case, therefore, they have been very convincing make-believes, because people believe in them; evident in the ways that these concepts have proven to be powerful motivators for human action. The term “action” here refers to the basic ways that a people perceive themselves relative to ‘others’, as well as explicit political actions staged against ‘others’. Therefore, it may be argued that, constructed or not, humans were predisposed for such contexts of identity and action. The social nature of humans may be a significant factor in this regard, and from which the distinction between “us” and “them” is inevitable.

\section{Identity as the Basis for Human Action}

There have been several cases of (collective) human actions throughout history. Take the example of social movements, for instance, of which there have been many examples, at various levels of identity, in this case ethnic, racial, and national levels. Horowitz\footnote{Horowitz, Donald L. \textit{Ethnic Groups in Conflict} (Los Angeles: University of California Press, 1985) 230-231}, for example, cites the ethnic clashes between Pakistan’s Baluch and Pathan groups, and Spain’s Catalans and Basque, among others. In terms of race, civil rights movements in the United States are perhaps the most telling cases, recently the Black Lives Matter (BLM) movement,\footnote{Jones, Leslie K. “\#BlackLivesMatter: An Analysis of the Movement as Social Drama.” \textit{Humanity \& Society vol.44}, No.1, 2020, pp.92-110} and most recently the ongoing Russian invasion of Ukraine.\footnote{Ukraine Crisis Lecture (2022)} The reason that these conflicts are identified with the levels assigned to them is evident; that they are conflicts between ethnicities, race groups, and nations. Underlying all of them is the question of identity, which has been in turn a central theme in social movements, as an example of collective human action.

As noted above, history is littered with several examples of social and political movements. Generally, in the language of Tarrow,\footnote{Tarrow, Sydney. \textit{Power in Movement: Social Movements and Contentious Politics}, Third Edition (Cambridge: Cambridge University Press, 2011) 16} these actions/movements are typically contentious, identified as uprisings against certain powers, so that they are a means of expressing common grievances. At the most basic level, therefore, identity has been a central drive in such actions, and underlying it are the notions of collective plight, therefore collective grievance against a given opposition. The problematic issue underlying these notions is: On what basis of identity should such actions be considered warranted? Particularly, the question is whether ethnicity, race, and nations are a valid basis for the collective identity that is perceived as necessary to inspire collective action.

\section{Ethnicity/Race/Nation as Constructed Concepts, But Also Distinct Manifestations of Problems that Validate Collective Social/Political Action}

Many would argue that ethnicity, race, and nationalities are facts. They cite several evidence to make their case. It is a fact that different populations around the world – and indeed within certain countries – speak different languages, which then become foundation for tribe and, therefore, tribal identity.\footnote{Chandra, Kanchan. “What Is Ethnic Identity and Does It Matter?” \textit{Annu. Rev. Polit. Sci. vol.9}, 2006, 397} Two, the concept arises from two elements: physical features (primarily), and cultural identity. Indeed, it seems inevitable to acknowledge the differences between various population groups in terms of their evident physical attributes, but culture proves much more problematic. As for nations, they are characterized in terms of countries, with borders that separate them from neighboring countries being the essential basis for that characterization: the scope of geographical reach and, in turn, the population that can be characterized as belonging to the same nation.\footnote{Ibid, p.397} However, looking far back and/or closely enough, the factuality of these concepts may simply be that they have been accepted as such, but that acceptance does not necessarily mean that they are true and, therefore, valid basis for human action, not even when that acceptance is widespread and collective. To use Brubaker’s\footnote{Brubaker, p.7} premise in \textit{Ethnicity without Groups,} such acceptance is based on commonsense, but commonsense cannot be relied upon here, particularly when the goal – if attainable – is to get to the intrinsic ‘truth’ of these phenomena.

Regardless, the basic premise here is that humans are social beings, which is arguably the basis of society – in this case ethnicity and nationality. In other words, humans have a tendency to live their lives associated with other lives. To Aristotle, for example, humans cannot live/exist in isolation from others, because they understand that their chances of survival improve when they cooperate with others.\footnote{Asikaogu, Joannes. “A Critical Review of Aristotle's view on Human sociality.” (2018) 37} In turn, this is how societies came to be, driven by the human urge to live in close proximity, where they would create identity relative to one another. Society may be what Brubaker refers to as “commonsense groupings.”\footnote{Brubaker, p.7} As it happens, societies define and distinguish themselves relative to others, and it may be from the basic notion of the ‘other’ – meaning ‘not us’ – that human groupings become a basis for, among others, cultural identity, kinship, ethnicity, and race.\footnote{Ibid} In other words, it becomes necessary for a grouping to find ways to distinguish itself from others, meaning whom to or not to consider kin. Groupings, therefore, become a basis for intra-group loyalty and, in turn, a basis for self-protection of the members against non-members. As noted above, collective human actions are driven by the perceptions of collective subjectivity to certain social/political patterns, meaning shared identity, particularly shared subjugation.

It so happens that, fictional or not, ethnicity, race, and nationality have become a basis for identity, and evidence points to several valid grievances along the lines. On this note, it is important to look at two important elements. One, the fundamental nature of societies (that is, their establishments and maintenance) as a basis for collective identity, which in turn lead to the collective emotional instincts that drive political action. Two, the political forces from outside that provoke the collective emotions that manifest in the form of political action.

\subsection{Ethnicity}

Ethnicity is often described on the basis of common language and culture. Language points to ‘tribes” as a group that share the same language, and culture manifests through day-to-day human behavior, based on certain commonly perceived and practices values and beliefs.\footnote{Chandra, pp.397-8} However, upon scrutiny, language and culture are themselves social constructions to begin with. Without going into the particular theories on the origins of language and the controversies they have caused, they all suggest that language was not always there, and that they emerged out of a need for social interaction.\footnote{Schwartz, Mila. \textit{Second Generation Immigrants: A Socio-Linguistic Approach of Linguistic Development Within the Framework of Family Language Policy}. In Mark Leikin, Mila Schwartz, \& Yishai Tobin (Eds.), \textit{Current Issues in Bilingualism: Cognitive and Socio-Linguistic Perspectives}, pp.119-135 (New York: Springer, 2012) 120} In other words, language is a transactional tool, and its adoption and/or usage is driven by the acquirer’s perception of its value. One could argue against this point, and yet the proof is obvious in the fact that second-generation immigrant children – and even more so the subsequent generations – will often take up the language of the new environment more than they will their first language. This point can be understood from the perspective of sociolinguistics, in which it is argued that the first language no longer has the same transactional significance as the new language, so the motivation for it is low.\footnote{Schwarts, p.120} The point made here is that no child is born to speak any one language, but that the environment will determine what language they speak and how they speak it. Therefore, as contextually-linked as it is, language fails as a basis for ethnic identity. Culture works the same way. As Brubaker\footnote{Brubaker, p.9} points out, culture is an aspect of essentialization and naturalization of people. Essentially, it means that one born in a given context is bombarded with the narrative and, indeed, examples from role models exhibiting what it means to belong to a given place, among a certain people.\footnote{Brubaker, p.9} Therefore, as Brubaker notes, ethnicity is constructed within a group, but also without; the latter emerging from outside talk: “policy analysis, media reports, and… constructivist academic writing [which] routinely frame accounts of ethnic, racial, and national conflict in groupist terms.”\footnote{Ibid} In that respect, it is possible for certain self-identifying ethnic groups to take action, to get their shared grievances addressed. The native tribes of America, for example, have a shared story – despite their given differences – of being the natives of the land, victimhood to colonial powers that stole land from them, and ultimately their placement in reserves. The smaller ethnic distinctions that go far back in history come to mean little in social/political action. At this point, that is the need for numbers, there is more focus on the larger similarities rather smaller differences – at least “small” in the face of the bigger ‘enemy’ they share.

\subsection{Race}

The concept of race is, perhaps, one of the most debated, and over time it has been rejected more intensely. Particularly questioned is the “scientific” or “genetic” notion of race that insists on defining a people as one or other in relation to various human abilities.\footnote{Bamshad, Michael, Stephen Wooding, Benjamin A. Salisbury, \& Clairbone Stephens J. “Deconstructing the Relationship between Genetics and Race.” \textit{Nature Reviews: Genetics vol. 5}, 2004, p.598} At its most extreme, there have been arguments, for instance, that Black people are genetically predisposed for certain health problems. However, while Black people do exhibit certain distinct health problems, W. E. DuBois attributed this fact to social inequality, not biological differences.\footnote{Yudell et al., p.564} Reiterating this point, Penner and Saperstein\footnote{Penner, Andrew M., \& Aliya Saperstein. “Engendering Racial Perceptions: An Intersectional Analysis of How Social States Shapes Race.” \textit{Gender \& Society vol.27}, No.3, 2013, pp.326} point out the problems with the genetic assumptions of race by highlighting the ways that social status have helped to shape the very notion. For example, because the vast majority of prisoners are likely to be Blacks and Latinos, there is a tendency to associate Blacks and Latinos with criminality. In other words, the social status of Blacks and Latinos subject them to unfair discriminatory treatment, and that in turn leads to certain racial notions regarding them, whether those notions are factual or not.

Therefore, one cannot deny the fact that certain groups of people face disproportionate degree of discrimination and various forms of human rights abuses. Minority groups in the Western world, for instance, have been found to face various forms of abuses and/or inequalities.\footnote{Brubaker, p.9} Indeed, civil rights movements has been focused on addressing these issues, and the BLM has continued that effort in recent years.\footnote{Jones, pp.92-110} Therefore, constructed or not, ethnic and racial identity in the US seems to have provided a basis for seeing and pointing out the fact that certain groups (minorities) are evidently more disadvantaged than others, and that disadvantage is intended by certain people. That seem valid basis for collective identity, which in turn drives collective action.

\subsection{Nation}

Perhaps of all three concepts, nation as a basis for human action is the most challenging, especially in certain circumstances. The Ukraine crisis\footnote{Ukraine Crisis Lecture (2022)} is only example of the difficulties in defining a nation, and distinguishing it from the so-called ethnicities within it. To the international community, the Ukraine crisis is one of nations, with Russia on the other side, and yet part of the conflict seems to stem from the fact that Ukraine is itself is an unrecognized state, among other states (like Transnistria [Moldova]) which were part of larger USSR.\footnote{Ibid} Looked at closely, then, the conflict raises questions as to whether the Ukraine crisis is an international war or an ethnic/inter-ethnic conflict. It is difficult to say.

The Ukraine crisis, therefore, points to the difficulty in defining nation, which in turn points to the historical problems that have characterized nation-building. Indeed, many nations grew out of colonialism, as European nations divided territories among themselves, and the process seems to have ignored ethnic identities that actually precede the creation of national borders. In fact, many national borders seem to have split ethnic groups into different nationalities. In Africa, for example, there are the so-called quasi-states, which (like Ukraine) are formally recognized and acknowledged as states on the international stage, but lack some key attributes of statehood.\footnote{Brubaker, Rogers \& David D. Laitin. “Ethnic and Nationalist Violence.” \textit{Annu. Rev. Sociol. Vol.24,} 1998, pp.423-452} Perhaps it comes as no surprise, therefore, that the Ukraine crisis remains confusing, reflecting the same confusions surrounding the dissolution of Yugoslavia.\footnote{Ibid, p.436}

As for collective action, some may argue that there are those who are being persecuted (such as in Ukraine) and who can then claim – as many Ukrainians have – nationhood as a basis for fighting back. However, as far as the definition of ‘nation’ go – that is, as unclear as it is – such claims are merely substitutes for a ‘mass of people’ facing common, adverse plight. The victimhood becomes the basis for collective identity, but the notion of nation is difficult – if at all possible – to find.

\section{Conclusion}

As the paper shows, the concepts of ethnicity, race, and nation are constructions of historical colonial and political processes. Ethnicity – despite its own conceptualization challenges – may be easier to characterize, but all pose challenges, particularly when used as a basis for oppression as has been the case throughout history. Ethnicity and race, for instance, have helped to see clearly the oppression of minority groups in the US, which then provide valid basis for collective identity – on the basis of shared experiences – and ultimately collective action. Nations are much harder to conceptualize, which is why the true nature of the Ukrainian crisis proves difficult to characterize. Still, there are a group of people facing shared grievances, and that does provide a basis for collective action.

\end{document}
